\chapter{Problem Statement}

\section{Rationale}\label{sec:rationale}

High-frequency bus operation is a sequential control problem: a holding or stop-skipping decision changes the headways, queues, loads, and downstream dwell times encountered by later buses. Fixed timetables cannot adapt after these quantities depart from their planned values \cite{Ceder2007,Daganzo2009}. Reinforcement learning is therefore relevant because it learns an action policy from repeated state transitions rather than using a fixed dispatch rule.

The empirical case is CapMetro Rapid Route 801 in Austin, Texas. CapMetro's public APC dataset covers July--December 2021 and contains event timestamps, route and direction codes, stop IDs, boarding and alighting counts, onboard load, dwell time, revenue travel time and distance, event coordinates, and quality indicators \cite{TexasCapMetroAPC2021}. Under identical cleaning rules, the reproducible audit found 455{,}654 clean stop events for Route 801 and 376{,}801 for Route 803. Route 801 also had more recorded boardings and more positive-time/positive-distance segments, which justifies its selection as the data-richer primary case. The decision does not imply superior or inferior service quality.

The one-direction study subset uses direction code 6 and contains 229{,}421 clean stop events, 184 service-day codes, and 29 distinct stop IDs. The direction is not assigned a compass label because a checksum-verified 2021 GTFS snapshot has not yet been acquired. Historical stop names, route shapes, scheduled headways, and capacity values therefore remain gated rather than being borrowed from the current schedule.

NOAA LCDv2 observations make an empirical ordinary-weather exposure feasible \cite{NOAALCDv2}. After the documented local-standard-time conversion, the nearest-observation join matched all primary APC events at both Camp Mabry and Austin-Bergstrom within 90 minutes, including 11{,}804 Camp Mabry rain-exposed events. Join coverage alone does not establish causality. Ordinary-rain effects must be estimated within comparable segment, time-of-day, and day-type strata; severe weather outside the observed support remains an explicitly synthetic stress test. Vehicle breakdowns are also synthetic because the APC dataset contains no failure-event field.

The resulting problem is to determine whether a parameter-sharing MARL controller, calibrated to observed CapMetro demand, dwell, load, and segment travel-time behavior, remains competitive with fixed-rule controllers when controlled demand surges, weather stress, and breakdowns are introduced without misrepresenting those synthetic variables as measured CapMetro incidents.

\section{Research Gap}

Existing Multi-Agent Reinforcement Learning frameworks for bus scheduling demonstrate measurable improvements in passenger waiting time and headway regularity, but these gains are validated almost exclusively in simulation environments that assume stationary passenger arrivals, symmetric Gaussian travel times, and a failure-free fleet \cite{Wang2020Holding,Rodriguez2023Cooperative,Wangsun}. The performance of MARL bus-scheduling models has not been characterized when these idealizing assumptions are simultaneously violated by (i) stochastic passenger demand surges, (ii) heavy-tailed weather-induced travel-time delays, and (iii) discrete bus breakdowns. Without this characterization, it cannot be determined whether reported MARL gains persist, degrade gracefully, or collapse under realistic operating disturbances, which in turn blocks the transition of MARL bus scheduling from simulation to real urban transit deployment.

The weather-disturbance class (W) in particular was arrived at through the literature survey conducted above: none of the MARL bus-scheduling studies reviewed models heavy-tailed, weather-induced travel-time delays, leaving the W column of the comparison unfilled. Its operational relevance is not assumed but established---rainfall materially reduces average speed and free-flow capacity on comparable corridors (Section~\ref{sec:background})---and it is made a controllable disturbance through the coefficient-of-variation-driven lognormal travel-time form validated by Patil et al.~\cite{Patil2025Conformal}. Weather therefore enters the study as a defined, parameterized stressor arrived at from the evidence rather than an untested label.

The CapMetro and NOAA sources narrow that gap in two ways. First, they replace assumed baseline demand, dwell, load, and segment travel-time profiles with reproducible public observations. Second, they permit observed ordinary-rain exposure to be separated from synthetic out-of-support stress. They do not supply passenger arrival times, capacity, breakdowns, or an authoritative historical schedule. The study therefore distinguishes measured calibration inputs, derived variables, simulated passenger states, and synthetic stressors in every claim.

\section{Research Objectives}

\subsection{Main Objective}

To evaluate a parameter-sharing Multi-Agent Reinforcement Learning controller for dynamic bus scheduling in a public-data-calibrated simulation of CapMetro Route 801, direction code 6, under the baseline and non-ideal operating conditions defined in Section~\ref{sec:scope}.

\subsection{Specific Objectives}
\begin{enumerate}
   \item \textbf{Environment Construction:} To construct a stochastic traffic simulation environment representing a defined traffic area.

   \textbf{Expected Output 1.1:} A simulation model of a defined operational scope, validated against empirical traffic data to establish an environment that accurately approximates real-world traffic.

   \textbf{Expected Output 1.2:} An environment simulator capable of introducing environmental variability through adjustable parameters to simulate non-ideal operating conditions.

   \item \textbf{Algorithm Implementation:} To develop and implement a Multi-Agent Reinforcement Learning (MARL) approach for dynamic bus scheduling.

   \textbf{Expected Output 2.1:} A formulated agent architecture defining individual bus units as agents with discrete state and action space with a continuous reward function that punishes irregularities and passenger service delays.

   \item \textbf{Performance Evaluation Under Ideal and Non-Ideal Conditions:} To evaluate the scheduling performance of the MARL approach under both ideal operating conditions and non-ideal conditions.

   \textbf{Expected Output 3.1:} A quantitative and statistical assessment of MARL scheduling performance under ideal conditions, measuring passenger waiting time, travel time, and headway coefficient of variation, tested against baseline methods.

   \textbf{Expected Output 3.2:} A quantitative and statistical assessment of MARL scheduling performance under non-ideal conditions, measuring passenger waiting time, travel time, and headway coefficient of variation, reported as degradation relative to ideal-conditions performance.
\end{enumerate}

\section{Significance of the Study}

This study contributes both practical and scientific significance, in that order: the operational question comes first, and the methodological contribution follows from how it is answered.

\textbf{Practical significance.} Route 801 recorded 810{,}309 boardings across 184 service days in the July--December 2021 window (Table~\ref{tab:capmetro-route-selection}). At that volume a small per-passenger improvement aggregates: a 30-second reduction in mean waiting time, applied across the recorded boardings, corresponds to roughly $6.8\times10^{3}$ passenger-hours over the same window. That figure illustrates the scale at which corridor control operates; it is not a performance target, and waiting time in this study is a simulated quantity rather than an observed one. The evaluation is further designed around the question an operator actually faces during a disruption, which is not average-case waiting time under quiet conditions but how far service degrades as conditions worsen---reported here across the baseline, combined-disturbance, and single-disturbance conditions. For CapMetro specifically, the route choice, data exclusions, NOAA join, and unresolved source gaps are documented closely enough that a third party can re-run or contest them.

\textbf{Scientific significance.} The study separates three layers that prior work commonly conflates: empirical ordinary operation, empirically joined weather exposure, and synthetic out-of-support stress. Keeping them distinct is what allows a robustness result to be attributed to the controller rather than to an undeclared modeling choice. The study also contributes the evaluation apparatus itself: a corridor model calibrated to public data, paired with configurable demand-surge, weather-stress, and breakdown generators, evaluated under matched random seeds across controllers. Because the corridor, the disturbance definitions, and the response variables are fixed while the underlying sources remain public, a later algorithm can be assessed on the same basis rather than on a private simulator, so performance claims in this area can be compared instead of reported in isolation.

\section{Scope and Limitations}\label{sec:scope}

\textbf{Scope.} The empirical scope is CapMetro Route 801, direction code 6, during July--December 2021. The simulation uses 29 observed stop IDs and a single simulated operating day per episode. The active fleet size, service horizon, scheduled headway, capacity, and control-stop set are implementation parameters that must be derived from verified sources or left to be finalized in the implementation phase. Each active bus is an agent that uses the same learned DDQN weights but produces its own action from a local observation.

The experimental activation matrix is fixed as follows. Baseline stochastic demand (D) and ordinary empirical segment travel-time variation (T) are present in training and every evaluation. Training domain-randomizes demand surge (S), weather stress (W), and breakdown (B). Stage A evaluates D+T only, with S/W/B disabled. Stage B evaluates D+T+S+W+B. Three ablations add S only, W only, or B only to the Stage A baseline. Thus W never replaces T, and an ablation labeled ``weather only'' means D+T+W rather than a deterministic corridor with weather as its sole source of randomness.

Ordinary observed rain is joined from Camp Mabry, with Austin-Bergstrom used for station sensitivity. Temperature, relative humidity, wind, and visibility may enter as control covariates, but only precipitation and rain/weather codes define the primary observed exposure. A lognormal coefficient-of-variation sweep may be used only for severe/extreme stress outside the empirical support and must be labeled synthetic. The breakdown process is likewise synthetic and must be reported separately from observed APC facts.

The MARL policy is compared with No Control, Forward Headway, and Even Headway using simulated passenger waiting time, bus travel time, and headway coefficient of variation over at least 30 matched-seed Monte Carlo runs per evaluation cell. Stage~A and Stage~B therefore describe evaluation conditions; training itself uses domain randomization across the disturbance classes throughout.

\textbf{Limitations.} (a) APC boardings do not identify passenger arrival times, so passenger waiting time is generated by the simulation and cannot be called directly observed. (b) APC coordinates are stop-event points, not continuous trajectories; calibration therefore uses segment travel time and event-level location quality rather than claiming continuous speed trajectories. (c) The units of \texttt{rev\_distance} depend on external odometer metadata and remain unresolved; travel-time RMSE is primary until units are verified. (d) Capacity and 2021 schedule semantics require separate authoritative sources. (e) Direction code 6 remains unlabeled until a compatible historical GTFS snapshot is checksum-verified. (f) No breakdown is inferred from missing or delayed APC records.

\textbf{Delimitations.} (g) Feeder-road traffic is not simulated explicitly; the agent observes only route-level headway, load, and queue states, so surrounding-traffic effects enter the model exclusively through the empirical demand and segment travel-time distributions already fit to the corridor. (h) The evaluation is based on repeated matched-seed Monte Carlo runs of a single simulated operating day rather than modeling long-term operational variation across multiple calendar days. (i) MARL control actions are applied only at the designated control stops selected under the study's stated criteria, and only demand surges, weather stress, and vehicle breakdowns are modeled as disturbance classes; other disturbance types (e.g., traffic accidents, driver behavior variability) are outside the study's scope. (j) The study is scoped to simulation only, with no on-road deployment or sim-to-real transfer claimed.